\documentclass[11pt,a4paper]{article}
\usepackage[margin=2.5cm]{geometry}
\usepackage{amsmath,amssymb,amsthm}
\usepackage{stmaryrd}
\usepackage{booktabs}
\usepackage{microtype}
\usepackage{xcolor}
\usepackage{listings}
\usepackage{parskip}
\usepackage[colorlinks=true,linkcolor=blue!50!black,urlcolor=blue!50!black]{hyperref}

\definecolor{cKey}{HTML}{1F4E8C}
\definecolor{cCom}{HTML}{6B7280}
\definecolor{cBg}{HTML}{F7F7F5}

\lstdefinestyle{pseudo}{
  basicstyle=\ttfamily\small,
  keywordstyle=\color{cKey}\bfseries,
  commentstyle=\color{cCom}\itshape,
  backgroundcolor=\color{cBg},
  frame=leftline, framerule=1.2pt, rulecolor=\color{cKey!30},
  xleftmargin=1.2em, framexleftmargin=1.2em,
  aboveskip=1em, belowskip=1em,
  columns=fullflexible, keepspaces=true,
  morecomment=[l]{--},
  morekeywords={type,data,where,let,in,if,then,else,case,of,for,while,
                return,function,and,or,not,do,forall,class,instance}
}
\lstset{style=pseudo}

\theoremstyle{definition}
\newtheorem{definition}{Definition}
\newtheorem{example}[definition]{Example}
\theoremstyle{plain}
\newtheorem{proposition}[definition]{Proposition}
\newtheorem{theorem}[definition]{Theorem}
\theoremstyle{remark}
\newtheorem*{remark}{Remark}

\newcommand{\dc}{\ensuremath{S \triangleleft P}}
\newcommand{\den}[1]{\ensuremath{\llbracket #1 \rrbracket}}

\title{A game on a directed container}
\author{Claude Opus 5 --- for Robin Langer}
\date{}

\begin{document}
\maketitle
\vspace{-2.5em}

\begin{abstract}
\noindent
The repository \texttt{stream} contains a fighting game built on the
\emph{stream comonad}. Stream is the degenerate directed container: one shape,
one future. This note defines the directing structure in general, gives
pseudocode, and reports what happened when a game was designed on a
non-degenerate one.

The short version: the game failed, twice, for reasons worth recording. The
mathematics it was built on survives intact. And the place where the director
\emph{does} earn its keep turned out not to be the new game at all, but the old
one --- the renormalisation that the shipped game now runs on is a statement
about the endomorphism monoid of a directed container.
\end{abstract}

\section{Containers}

A \emph{container} is a pair \dc{} consisting of a set $S$ of \emph{shapes} and,
for each shape $s : S$, a set $P\,s$ of \emph{positions} in that shape. Shapes
are templates; positions are the blanks in a template that data can be poured
into. A container denotes a functor on $\mathbf{Set}$:
\[
  \den{\dc}\,X \;=\; \sum_{s : S} X^{P\,s},
\]
an element of which is a shape together with a labelling of its positions. A
list is the container with $S = \mathbb{N}$ (the length) and
$P\,s = \{0,\dots,s-1\}$. A stream is the container with $S = 1$ and
$P\,\ast = \mathbb{N}$.

Containers denote functors, and functors alone. To get a \emph{comonad} --- to
get \texttt{extract} and \texttt{duplicate}, and therefore \texttt{extend} ---
something more is needed, and that something is what a director supplies.

\section{The director}

\begin{definition}[Ahman--Chapman--Uustalu]\label{def:dc}
A \emph{directed container} is a container \dc{} together with three operations
\[
\begin{array}{ll}
  \downarrow \;:\; \Pi\,s : S.\; P\,s \to S
    & \text{the \emph{subshape} at a position,}\\[2pt]
  o \;:\; \Pi\,\{s : S\}.\; P\,s
    & \text{the \emph{root} position,}\\[2pt]
  \oplus \;:\; \Pi\,\{s : S\}.\;\Pi\,p : P\,s.\; P\,(s \downarrow p) \to P\,s
    & \text{\emph{translation} into a global position,}
\end{array}
\]
satisfying two shape equations and three position equations:
\begin{align}
  s \downarrow o &= s \\
  s \downarrow (p \oplus p') &= (s \downarrow p) \downarrow p' \\
  p \oplus o &= p \\
  o \oplus p &= p \\
  (p \oplus p') \oplus p'' &= p \oplus (p' \oplus p'')
\end{align}
\end{definition}

I will call the triple $(\downarrow, o, \oplus)$ the \emph{director} of the
container, and $\downarrow$ itself the \emph{directing map}: given where you are
and a position you can see from there, it \emph{directs} you to the world as
seen from that position.

\begin{remark}
This is my naming, not the literature's. Ahman, Chapman and Uustalu call
$\downarrow$ the subshape operation, $o$ the root and $\oplus$ translation, and
do not use the word ``director'' anywhere in the paper. The polynomial-functor
school does use \emph{directions} for what containers call positions, which is
the nearest standard usage. If ``director'' means something specific to you or
to Kodamai that differs from Definition~\ref{def:dc}, this is the one place in
the document that needs changing.
\end{remark}

Read the laws as geometry. Law (1) says the subshape at the root is the whole
shape --- looking from where you already are changes nothing. Law (3) says the
root of the subshape at $p$ translates back to $p$ itself, so the subshape given
by a position really is \emph{rooted} at that position. Law (2) says that going
to $p$ and then to $p'$ lands on the same shape as going directly to
$p \oplus p'$. Laws (3)--(5) are the monoid laws, except that the three
positions live in three different sets, so it is a monoid only in a dependent
sense. Laws (1)--(2) are the laws of a monoid \emph{action}.

That last observation is the whole content of the definition. A directed
container is a monoid-like structure of positions acting on shapes, where the
monoid you get depends on where you are standing.

\subsection{How the director is used: it is exactly a comonad}

\begin{theorem}[ACU]
$\den{\dc}$ carries a comonad structure precisely when \dc{} carries a director.
\end{theorem}

The comonad is built from the director directly:
\[
\begin{array}{ll}
  \varepsilon\,(s, f) &=\; f\,(o) \\[3pt]
  \delta\,(s, f) &=\; \bigl(s,\;\lambda p.\;
      (\,s \downarrow p,\;\; \lambda p'.\; f\,(p \oplus p')\,)\bigr)
\end{array}
\]
$\varepsilon$ reads the label at the root: \emph{what is true here}.
$\delta$ replaces the label at each position $p$ by the entire structure
re-rooted at $p$: \emph{the world as it looks from there}. The type-checking of
$\delta$ is exactly laws (1)--(2), and its comonad laws are exactly (3)--(5).
And then
\[
  \mathrm{extend}\; :\; (\den{\dc}A \to B) \;\to\; \den{\dc}A \to \den{\dc}B,
  \qquad \mathrm{extend}\,g \;=\; \den{\dc}g \circ \delta,
\]
which runs a rule that needs to see a whole neighbourhood at \emph{every}
position at once. That is the operation \texttt{fmap} cannot express, and it is
the reason the structure is worth having.

\subsection{Stream is the degenerate case}

Take $S = 1$, $P\,\ast = \mathbb{N}$, $s \downarrow p = s$, $o = 0$,
$p \oplus p' = p + p'$. Laws (1)--(2) are trivial because the shape never
changes; laws (3)--(5) are $p+0=p$, $0+p=p$ and associativity of $+$. Then
$\den{\dc}X = X^{\mathbb{N}}$ is streams, $\varepsilon$ is \texttt{head}, and
$\delta$ is \texttt{tails}. The file \texttt{StreamComonad.hs} prints $\delta$
applied to $0,1,2,\dots$ and observes that it is the addition table. That is
law (2) with $\downarrow$ erased.

\begin{center}
\begin{tabular}{lll}
\toprule
 & stream & general directed container \\
\midrule
shapes $S$        & one                 & many: the world depends on where you stand \\
$s \downarrow p$  & $s$, unchanged      & a \emph{different} world \\
$\oplus$          & $+$, commutative    & composition, \textbf{not} commutative \\
futures           & exactly one         & \textbf{you choose} \\
\bottomrule
\end{tabular}
\end{center}

\subsection{The director we use: a graph}

Let $G$ be a directed graph. Define
\[
  S = \text{vertices}, \quad
  P\,v = \{\text{finite paths out of } v\}, \quad
  v \downarrow p = \mathrm{target}(p), \quad
  o = \varepsilon, \quad
  p \oplus p' = p \cdot p'.
\]
All five laws hold, and this is the free category on $G$. Now $\oplus$ is
genuinely non-commutative and $\downarrow$ genuinely moves you.

This is not a coincidence of the example. Ahman and Uustalu proved that
\emph{directed containers are precisely small categories}: objects are shapes,
morphisms out of $s$ are positions in $P\,s$, $\downarrow$ is codomain, $o$ is
the identity and $\oplus$ is composition. A director is a category wearing the
notation of a data structure. Section~\ref{sec:endo} is where that fact does
some work.

\section{The world: a Rauzy graph}

Fix an irrational slope $\alpha$ and let $w$ be the Sturmian word of slope
$\alpha$. Its \emph{Rauzy graph of order $m$} has one vertex per factor of
length $m$ --- a \emph{room} --- and an edge $u \to v$ whenever $u = a\sigma$
and $v = \sigma b$ with $a\sigma b$ a factor, which \emph{emits} the bit $b$.

Walking the graph emits a bit sequence, and \textbf{the room you stand in
\emph{is} the last $m$ bits you emitted}. The rooms are not extra state; they
are your own recent history.

\begin{proposition}\label{prop:necklace}
The Rauzy graph of order $m$ of a Sturmian word has exactly $m+1$ rooms and
$m+2$ doors, and exactly one room has two doors, whose bits differ.
\end{proposition}

Checked in \texttt{Rauzy.hs}, 40 slopes $\times$ 6 orders, $240/240$ on each
count. The $m+1$ is the Sturmian complexity function $p(m) = m+1$: the ``one
split at every window length'' of \texttt{splitting.pdf} is \textbf{one
branching room}. Depending on whether the right- and left-special factors
coincide, the graph is a \emph{theta} (two arcs reconverging) or a
\emph{figure-eight} (two loops meeting at a point).

\section{The game that was designed, and why it failed}

The design: you are inside the schedule rather than watching it go past. Each
beat you choose a door; the bit it emits is the monster's action that beat.
Combat is unchanged from \texttt{Game.hs}. Most rooms have one door --- no
choice, exactly the old game. One room has two.

It failed twice, and both failures are worth more than the design was.

\subsection{The fiction would not go into words}

Asked how a monster can have doors, the answer is that it cannot. ``Rooms'' and
``doors'' were imported from a \emph{picture} of the graph without checking that
the metaphor survived the transplant. The vertices are not places; they are
\emph{situations} --- the last $m$ beats of the exchange --- and the edges are a
positioning choice that constrains what the monster can do next.

\subsection{Stated plainly, the mechanic died}

Once the fiction is stated properly --- ``you choose whether the monster strikes
you'' --- the next question is immediate: why would you ever choose to be
struck? The intended answer was a trade: the free hit now versus the density of
the loop it commits you to. Measured over 400 slopes at each of five orders:

\begin{center}
\begin{tabular}{lr}
\toprule
does taking the beat that is open \emph{now} also give the denser loop? & \\
\midrule
$m = 3,\,5,\,8,\,13,\,21$ & $1993/1993$ \\
\bottomrule
\end{tabular}
\end{center}

\textbf{They never disagree.} The greedy step is always also the better
long-run loop, so the one decision in the world has a dominant answer requiring
no lookahead at all. What remained was an information problem --- find out which
door is open --- with no strategy in it.

\begin{remark}
An earlier draft of this document called the reference player ``Karp'' and
described the branch as a choice between convergents. Both were wrong, and the
dominance result above is why. Plain greed is optimal; the max-mean-cycle
framing was decorative. That claim is withdrawn.
\end{remark}

\section{Pseudocode}

\subsection{The director, as an interface}

\begin{lstlisting}
type Shape
type Position         -- indexed by a shape: Position s

class DirectedContainer where
  down  :: Shape -> Position s -> Shape          -- s |> p     the directing map
  root  :: Position s                            -- o
  plus  :: Position s -> Position (down s p) -> Position s   -- p <+> p'

-- LAWS
--   down s root          == s            plus p root          == p
--   down s (plus p p')   == down (down s p) p'
--   plus root p          == p
--   plus (plus p p') p'' == plus p (plus p' p'')
\end{lstlisting}

\subsection{The comonad it generates}

\begin{lstlisting}
data Den x = Den Shape (Position -> x)      -- [| S <| P |] x

extract :: Den x -> x
extract (Den s f) = f root                  -- what is true HERE

duplicate :: Den x -> Den (Den x)
duplicate (Den s f) =
  Den s (\p -> Den (down s p)               -- the world re-rooted at p ...
                   (\p' -> f (plus p p')))  -- ... relabelled by translation

extend :: (Den a -> b) -> Den a -> Den b
extend g = fmap g . duplicate               -- run g at EVERY position at once
\end{lstlisting}

\noindent
Specialising to $S = 1$, $\texttt{down}\,s\,p = s$, $\texttt{root} = 0$,
$\texttt{plus} = (+)$ turns \texttt{duplicate} into \texttt{tails} and
$\texttt{plus}\,p\,p'$ into $p + p'$ --- the addition table.

\subsection{Building the world, and the fight}

\begin{lstlisting}
sturmian alpha rho =
  [ floor ((k+1)*alpha + rho) - floor (k*alpha + rho) | k <- [0..] ]

rauzy m w = groupByRoom [ (window m i, (w !! (i+m), window m (i+1)))
                        | i <- [0 .. prefix] ]
  where window m i = take m (drop i w)

fight graph start playerHP monsterHP policy =
  loop start playerHP monsterHP emptyMap
  where
    loop room php mhp known
      | mhp <= 0 = WIN
      | php <= 0 = DEAD
      | otherwise =
          let doors       = doorsOf graph room
              (i, attack) = policy room known     -- the coKleisli arrow
              (bit, next) = doors !! i            -- <-- this IS  room |> p
              php'        = if attack and bit == STRIKE then php - 5 else php
              mhp'        = if attack and bit == OPEN   then mhp - 2 else mhp
          in loop next php' mhp' (insert (room, i) bit known)
\end{lstlisting}

\noindent
\texttt{doors !!\ i} is the directing map applied to a length-one position, and
\texttt{known} is a \emph{partial} element of $\den{\dc}\,\texttt{Bit}$: a
partial function on positions, defined exactly on the doors already walked.
Scouting enlarges its domain.

\section{What was measured}

All figures 300 fights per cell, $\alpha$ uniform in $[0.25, 0.75]$.

\subsection{The mathematics that survives}

\begin{proposition}\label{prop:conv}
The open-densities of the two cycles through the branching room are the best
rational approximations to $1 - \alpha$ from above and from below, among
fractions with denominator at most $m+1$.
\end{proposition}

Agreement: $300/300$ at $m = 3, 5, 8, 13$ and $299/300$ at $m = 21$. Worked
cases at $m = 8$:

\begin{center}
\begin{tabular}{llll}
\toprule
$\alpha$ & $1-\alpha$ & cycle densities & best approximations \\
\midrule
0.3026 & 0.6974 & $5/7,\ 2/3$ & $5/7$ and $2/3$ \\
0.3490 & 0.6510 & $5/8,\ 2/3$ & $2/3$ and $5/8$ \\
0.6536 & 0.3464 & $1/3,\ 3/8$ & $3/8$ and $1/3$ \\
0.7000 & 0.3000 & $1/3,\ 2/7$ & $1/3$ and $2/7$ \\
\bottomrule
\end{tabular}
\end{center}

\begin{remark}[a correction]
I first stated this with denominator at most $m$, and it failed on roughly 30\%
of slopes. Every failure had denominator exactly $m+1$ --- unsurprising, since
the graph has $m+1$ rooms and no simple cycle can be longer. With the corrected
bound the agreement is $1499/1500$.
\end{remark}

\subsection{The design that does not}

Even before the dominance result, the decision was worth almost nothing at any
usable map size. Mean damage over 120 beats with no HP pool, so that nothing
hides behind a win threshold:

\begin{center}
\begin{tabular}{rrrrrr}
\toprule
$m$ & first arc & scouts both & oracle & arc gap $\Delta d$ & doors differ by \\
\midrule
3  & 104.1 & 131.3 & 137.0 & 0.1389 & 33.3 dmg \\
5  & 112.2 & 123.0 & 130.7 & 0.0854 & 20.5 dmg \\
8  & 114.8 & 114.5 & 125.0 & 0.0398 & \phantom{0}9.6 dmg \\
13 & 109.0 & 107.4 & 122.7 & 0.0200 & \phantom{0}4.8 dmg \\
21 & 103.7 & \phantom{0}98.1 & 121.4 & 0.0095 & \phantom{0}2.3 dmg \\
\bottomrule
\end{tabular}
\end{center}

The gap decays roughly as $m^{-1.4}$, because the two return times are typically
of very unequal size. At $m = 8$ --- the size I had planned to ship --- the two
doors differ by under 10 damage in a fight of 120.

\subsection{And pursuit would not have saved it}

A chase (you move two steps, the monster one) needs a board with somewhere to
hide. Sturmian Rauzy graphs have exactly one fork at every order, by
Proposition~\ref{prop:necklace}; and Arnoux--Rauzy words, whose complexity is
$2m+1$, also have exactly one, merely of out-degree three:

\begin{center}
\begin{tabular}{lrrrr}
\toprule
 & $m=5$ & $m=8$ & $m=13$ & $m=21$ \\
\midrule
Sturmian, branch vertices   & 1 & 1 & 1 & 1 \\
\quad out-degrees           & \multicolumn{4}{l}{one vertex of degree 2, rest degree 1} \\
Tribonacci, branch vertices & 1 & 1 & 1 & 1 \\
\quad out-degrees           & \multicolumn{4}{l}{one vertex of degree 3, rest degree 1} \\
\bottomrule
\end{tabular}
\end{center}

One fork per lap is one binary guess per lap, which is the decision structure
the original game already has. Tribonacci's three-way fork is the version that
could not be straddled, and is where I would look if pursuit were wanted.

\section{Where the director earns its keep}\label{sec:endo}

None of the above uses the director for anything a bare graph could not do: only
length-one positions, a $\downarrow$ with no shape content, and
\texttt{extend} never called once. But the identification
\emph{directed container $=$ small category} has a consequence that is not
decorative.

In a small category, the morphisms from an object to itself form a
\textbf{monoid}. Translated back: the positions $p \in P\,s$ with
$s \downarrow p = s$ are closed under $\oplus$, contain $o$, and are associative
--- by laws (1)--(5) and nothing else. For the Rauzy graph, the loops at the
branching room are freely generated by the two simple cycles, because every loop
there decomposes uniquely into them.

So a trajectory through the world, read \emph{in the endomorphism monoid of the
branching object}, is an infinite word on two letters: which loop was taken,
each time round.

\begin{proposition}
That word is Sturmian again.
\end{proposition}

Checked at $m = 2, 3, 5, 8, 13$, six slopes each: complexity $k+1$ in all
$30/30$ cases.

This is the classical renormalisation of a Sturmian word --- passing to the
derived word shifts the slope by the Gauss map, i.e.\ consumes one partial
quotient of $\alpha$ --- but stated in the container's own language it says
something sharper: \textbf{the directed container contains a copy of itself,
sitting in the endomorphism monoid of its one branching object.}

\subsection{And this is what the shipped game now runs on}

At $m = 2$ the branching room is the factor \texttt{.H}, and the loop taken is
just the next beat. Robin's reading of the game is then a set of theorems rather
than heuristics: for $1/2 < \alpha < 2/3$ the word contains no \texttt{..} and
no \texttt{HHH} ($0/300$ each), so of the three length-2 contexts exactly one is
undetermined:
\[
  \texttt{H.} \to \texttt{H} \ \text{certain},
  \qquad
  \texttt{HH} \to \texttt{.} \ \text{certain},
  \qquad
  \texttt{.H} \to \texttt{?} \ \text{the split.}
\]
Together with his replacement of the fight clock by \emph{chip damage} --- a
blocked strike costs 1, so stalling is paid for in the same currency as
everything else --- this gives a game with no dominant strategy and a ceiling a
table cannot reach. 300 fights, $\alpha \in [0.52, 0.62]$, 60 HP each:

\begin{center}
\begin{tabular}{lr}
\toprule
strategy & wins \\
\midrule
attack every beat & \phantom{00}0\% \\
swing only after \texttt{HH} --- never wrong, too slow to outrun the chip & \phantom{00}0\% \\
swing after any strike --- takes the split blind & \phantom{0}23\% \\
memory-8 table on the raw beats (256 contexts) & \phantom{0}82\% \\
\textbf{renormalising, depth 4} (16 contexts on the derived word) & \textbf{92\%} \\
perfect prediction & 100\% \\
\bottomrule
\end{tabular}
\end{center}

A 16-context recursive player beats a 4096-context memoriser, because it learns
$\alpha$'s next partial quotient instead of memorising more of the word. The
memory-8 table flatlines at 82--83\% no matter how deep it goes; the
renormaliser does not, and Section~\ref{sec:endo} is why.

\section{What a richer game would have to use}

The test I should have applied at the start: \textbf{if the mechanic still works
after you delete the director, the director was decoration.} Three parts of
Definition~\ref{def:dc} went unused, and each is a mechanic.

\begin{enumerate}\itemsep4pt
\item \textbf{$P\,s$ is all paths, not all steps.} Make a move a declared path
      $p \in P\,s$: commit $k$ beats, blind for $k$, land at $s \downarrow p$.
      Law (5) is exactly what makes a combo mean one thing however you bracket
      it, and non-commutativity of $\oplus$ is the mathematical content of
      \emph{initiative} --- invisible in the stream case, where positions
      commute.
\item \textbf{$\downarrow$ can consume the world.} The worst thing in the
      current design is \texttt{beats = 120}, a constant with no mathematical
      content. In the non-empty-list director, $s \downarrow p = s - p$: moving
      does not relocate you, it \emph{spends} the world. Then ``out of time'' is
      ``$s \downarrow p$ has no positions left'', and irreversibility is felt
      rather than asserted.
\item \textbf{\texttt{extend} is the verb no functor has.} Make it a
      consumable: spend a \emph{read}, and a fixed rule runs over your whole
      forward cone at once. A game whose central action is \texttt{extend} does
      not typecheck without $\delta$, and $\delta$ does not exist without laws
      (1)--(2). That is a mechanic the director cannot be deleted from.
\end{enumerate}

\noindent
And morphisms give the difficulty dial for free: ``drop the last letter'' is a
functor $G_{m+1} \to G_m$, so bounded memory is a \emph{quotient} directed
container and levelling up is climbing one rung of a tower whose inverse limit
is the Sturmian system itself.

\section{Status}

In the repository and verified: \texttt{Rauzy.hs}
(Proposition~\ref{prop:necklace}, $240/240$), \texttt{bench/Warren.hs}
(Proposition~\ref{prop:conv}, $1499/1500$), \texttt{bench/Chip.hs} and
\texttt{Game.hs} (the chip scoring and the renormalising player). Designed but
not written: \texttt{Directed.hs}, checking the five laws exhaustively on a
container where $\oplus$ is non-commutative.

Not built, deliberately: the Warren. Its mathematics is above; its mechanic has
a dominant answer. The next thing actually worth measuring is a renormaliser of
depth $d > 1$, recursing all the way down the continued fraction of $\alpha$
rather than stopping after one level.

\vspace{1em}\noindent\textbf{Sources.}
D.~Ahman, J.~Chapman, T.~Uustalu, \emph{When is a container a comonad?}, LMCS
10(3:14), 2014, \url{https://lmcs.episciences.org/894/pdf} ---
Definition~\ref{def:dc} is their Definition 3.1.
D.~Ahman, T.~Uustalu, \emph{Directed containers as categories}, MSFP 2016,
\url{https://arxiv.org/pdf/1604.01187}.

\end{document}
